\documentclass[11pt]{article}

% ============================================================
% Packages
% ============================================================
\usepackage[margin=1in]{geometry}
\usepackage{amsmath,amssymb,mathtools,amsthm}
\usepackage{enumitem}
\usepackage{hyperref}
\usepackage{xcolor}
\usepackage{microtype}

\theoremstyle{plain}
\newtheorem{theorem}{Theorem}[section]
% ============================================================
% Macros
% ============================================================
\newcommand{\R}{\mathbb{R}}
\newcommand{\Pbb}{\mathbb{P}}
\newcommand{\Qbb}{\mathbb{Q}}
\newcommand{\E}{\mathbb{E}}
\newcommand{\dd}{\,\mathrm{d}}
\newcommand{\F}{\mathcal{F}}

\newcommand{\1}{\mathbf{1}}
\newcommand{\Var}{\mathrm{Var}}
\newcommand{\Cov}{\mathrm{Cov}}

% ============================================================
% Environments
% ============================================================
\theoremstyle{definition}
\newtheorem{exercise}{Exercise}[section]

\newenvironment{solution}
{\par\medskip\noindent\textbf{Solution.}\par\medskip}
{\hfill$\square$\par\medskip}

\newenvironment{solutionpart}[1]
{\par\medskip\noindent\textbf{#1.}\par\smallskip}
{\par\medskip}

% ============================================================
% Title
% ============================================================
\title{Option Pricing --- Correction Sheet}
\author{Nathan Sauldubois}
\date{MSFE6233 --- Refresher}

\begin{document}

\maketitle
\vspace{-1em}

\noindent\textbf{Lecture title.}
\emph{Brownian Motion, Stochastic Calculus, Girsanov Theorem, Option Pricing and Delta Hedging}

\vspace{0.5em}
\hrule
\vspace{1em}

\tableofcontents
\bigskip

% ============================================================
\section{Brownian Motion}
% ============================================================
\begin{solution} \noindent 
\textbf{Compute $\int_0^t e^{B_s}\,\mathrm{d} B_s$ without $\mathrm{d} B$}

\noindent Let
\[
f(x)=e^x.
\]
Then
\[
f'(x)=e^x,
\qquad
f''(x)=e^x.
\]
By It\^o's formula,
\[
\dd e^{B_t}
=
e^{B_t}\,\dd B_t
+
\frac12 e^{B_t}\,\dd t.
\]
Integrating from $0$ to $t$, we get
\[
e^{B_t}-e^{B_0}
=
\int_0^t e^{B_s}\,\dd B_s
+
\frac12\int_0^t e^{B_s}\,\dd s.
\]
Since $B_0=0$, we have $e^{B_0}=1$. Therefore,
\[
e^{B_t}-1
=
\int_0^t e^{B_s}\,\dd B_s
+
\frac12\int_0^t e^{B_s}\,\dd s.
\]
Hence
\[
\boxed{
\int_0^t e^{B_s}\,\dd B_s
=
e^{B_t}-1
-
\frac12\int_0^t e^{B_s}\,\dd s.
}
\]
\end{solution}

% -----------------------------------------------------------

\begin{solution}
\textbf{Compute $e^{\int_0^t f(s)\,\mathrm{d} B_s}$ without $\mathrm{d} B$}
Assume that $f:[0,T]\to\mathbb{R}$ is deterministic and satisfies
\[
\int_0^T f(s)^2\,\dd s<\infty.
\]
Define
\[
X_t:=\int_0^t f(s)\,\dd B_s.
\]
Then $X$ is a continuous martingale with quadratic variation
\[
\langle X\rangle_t
=
\int_0^t f(s)^2\,\dd s.
\]

We want to compute
\[
e^{\int_0^t f(s)\,\dd B_s}
=
e^{X_t}
\]
in a form where the stochastic integral is isolated or expressed using an exponential martingale.

By It\^o's formula applied to the function $x\mapsto e^x$, we get
\[
\dd e^{X_t}
=
e^{X_t}\,\dd X_t
+
\frac12 e^{X_t}\,\dd\langle X\rangle_t.
\]
Since
\[
\dd X_t=f(t)\,\dd B_t
\qquad\text{and}\qquad
\dd\langle X\rangle_t=f(t)^2\,\dd t,
\]
we obtain
\[
\dd e^{X_t}
=
e^{X_t}f(t)\,\dd B_t
+
\frac12 e^{X_t}f(t)^2\,\dd t.
\]
Therefore,
\[
e^{X_t}
=
1
+
\int_0^t e^{X_s}f(s)\,\dd B_s
+
\frac12\int_0^t e^{X_s}f(s)^2\,\dd s.
\]

This identity still contains a stochastic integral. The standard way to remove the drift term is to introduce the exponential martingale
\[
M_t
:=
\exp\left(
\int_0^t f(s)\,\dd B_s
-
\frac12\int_0^t f(s)^2\,\dd s
\right).
\]
Equivalently,
\[
M_t
=
e^{X_t-\frac12\langle X\rangle_t}.
\]

Let us verify that $M$ is indeed a local martingale. Applying It\^o's formula to
\[
M_t=\exp\left(X_t-\frac12\langle X\rangle_t\right),
\]
we get
\[
\dd M_t
=
M_t\,\dd X_t
=
M_t f(t)\,\dd B_t.
\]
Hence
\[
M_t
=
1+\int_0^t M_s f(s)\,\dd B_s.
\]
In particular, $M$ has no drift term. Under standard integrability assumptions, for example Novikov's condition, $M$ is a true martingale. Since $f$ is deterministic and square-integrable on $[0,T]$, this is satisfied.

Thus,
\[
\boxed{
\exp\left(
\int_0^t f(s)\,\dd B_s
-
\frac12\int_0^t f(s)^2\,\dd s
\right)
}
\]
is the exponential martingale associated with $\int_0^t f(s)\,\dd B_s$.

Finally, rearranging the definition of $M_t$, we obtain
\[
\boxed{
e^{\int_0^t f(s)\,\dd B_s}
=
\exp\left(
\frac12\int_0^t f(s)^2\,\dd s
\right)
M_t.
}
\]
Equivalently,
\[
\boxed{
e^{\int_0^t f(s)\,\dd B_s}
=
\exp\left(
\frac12\int_0^t f(s)^2\,\dd s
\right)
\exp\left(
\int_0^t f(s)\,\dd B_s
-
\frac12\int_0^t f(s)^2\,\dd s
\right).
}
\]

In particular,
\[
\int_0^t f(s)\,\dd B_s
\sim
\mathcal{N}\left(
0,\int_0^t f(s)^2\,\dd s
\right),
\]
and therefore
\[
\E\left[
e^{\int_0^t f(s)\,\dd B_s}
\right]
=
\exp\left(
\frac12\int_0^t f(s)^2\,\dd s
\right).
\]
\end{solution}

% ============================================================
\section{Girsanov}
% ============================================================


\begin{solution}\textbf{Girsanov and the inverse Gaussian law}
We consider
\[
X_t=\mu t+\sigma B_t,
\qquad
\mu>0,\quad \sigma>0,
\]
and
\[
\tau_a:=\inf\{t\ge 0:X_t=a\},
\qquad a>0.
\]

\begin{solutionpart}{Question (a)}
We show that
\[
\tau_a<\infty
\qquad \mathbb{P}\text{-a.s.}
\]

Since
\[
X_t=\mu t+\sigma B_t,
\]
we have
\[
\frac{X_t}{t}
=
\mu+\sigma\frac{B_t}{t}.
\]
It is a standard property of Brownian motion that
\[
\frac{B_t}{t}\longrightarrow 0
\qquad \mathbb{P}\text{-a.s.}
\]
as $t\to\infty$. Hence
\[
\frac{X_t}{t}\longrightarrow \mu>0
\qquad \mathbb{P}\text{-a.s.}
\]
Therefore,
\[
X_t\longrightarrow +\infty
\qquad \mathbb{P}\text{-a.s.}
\]
Since $X$ has continuous paths and $X_0=0<a$, the intermediate value theorem implies that $X$ must hit the level $a$ in finite time. Thus
\[
\tau_a<\infty
\qquad \mathbb{P}\text{-a.s.}
\]
\end{solutionpart}

\begin{solutionpart}{Question (b)}
Define
\[
\theta:=\frac{\mu}{\sigma},
\qquad
Z_t:=\exp\left(-\theta B_t-\frac12\theta^2t\right).
\]
This is the exponential martingale associated with the Brownian integral
\[
-\int_0^t \theta\,\dd B_s.
\]
Since $\theta$ is constant, Novikov's condition is satisfied:
\[
\E^\mathbb{P}\left[
\exp\left(\frac12\int_0^T \theta^2\,\dd s\right)
\right]
=
\exp\left(\frac12\theta^2T\right)
<\infty.
\]
Therefore, $(Z_t)_{0\le t\le T}$ is a true martingale for every $T>0$.

We recall Girsanov Theorem : 

\begin{theorem}[Girsanov theorem with Novikov's condition]
Let $(\Omega,\mathcal F,(\mathcal F_t)_{0\le t\le T},\mathbb P)$ be a filtered probability space satisfying the usual conditions, and let
$B=(B_t)_{0\le t\le T}$ be a $d$-dimensional Brownian motion under $\mathbb P$.

Let $\theta=(\theta_t)_{0\le t\le T}$ be an $\mathbb R^d$-valued progressively measurable process such that
\[
\int_0^T |\theta_s|^2\,\mathrm ds<\infty,
\qquad \mathbb P\text{-a.s.}
\]
Assume moreover that Novikov's condition holds:
\[
\mathbb E^{\mathbb P}\left[
\exp\left(
\frac12\int_0^T |\theta_s|^2\,\mathrm ds
\right)
\right]<\infty.
\]
Define the stochastic exponential
\[
Z_t
:=
\exp\left(
-\int_0^t \theta_s\cdot \mathrm dB_s
-\frac12\int_0^t |\theta_s|^2\,\mathrm ds
\right),
\qquad 0\le t\le T.
\]
Then $(Z_t)_{0\le t\le T}$ is a true $\mathbb P$-martingale, and one may define a probability measure $\mathbb Q$ on $\mathcal F_T$ by
\[
\frac{\mathrm d\mathbb Q}{\mathrm d\mathbb P}
=
Z_T.
\]
Under $\mathbb Q$, the process
\[
\widetilde B_t
:=
B_t+\int_0^t \theta_s\,\mathrm ds,
\qquad 0\le t\le T,
\]
is a $d$-dimensional Brownian motion.
\end{theorem}

We define a new probability measure $\mathbb{Q}$ on $\mathcal{F}_t$ by
\[
\frac{\dd\mathbb{Q}}{\dd\mathbb{P}}\Big|_{\mathcal{F}_t}
=
Z_t.
\]
By Girsanov's theorem, the process
\[
\widetilde B_t
:=
B_t+\theta t
\]
is a standard Brownian motion under $\mathbb{Q}$.
\end{solutionpart}

\begin{solutionpart}{Question (c)}
We rewrite $X_t$ in terms of $\widetilde B_t$.
Since
\[
\widetilde B_t=B_t+\theta t,
\]
we have
\[
B_t=\widetilde B_t-\theta t.
\]
Therefore,
\[
X_t
=
\mu t+\sigma B_t
=
\mu t+\sigma(\widetilde B_t-\theta t).
\]
Using $\theta=\mu/\sigma$, this becomes
\[
X_t
=
\mu t+\sigma\widetilde B_t-\sigma\frac{\mu}{\sigma}t
=
\sigma\widetilde B_t.
\]
Hence, under $\mathbb{Q}$,
\[
X_t=\sigma\widetilde B_t.
\]

Therefore,
\[
\tau_a
=
\inf\{t\ge0:X_t=a\}
=
\inf\{t\ge0:\sigma\widetilde B_t=a\}
=
\inf\{t\ge0:\widetilde B_t=a/\sigma\}.
\]
Thus, under $\mathbb{Q}$, $\tau_a$ is the first hitting time of the level
\[
b:=\frac{a}{\sigma}
\]
by a standard Brownian motion.
\end{solutionpart}

\begin{solutionpart}{Question (d)}

\begin{theorem}[Reflection principle]
Let $(B_t)_{t\ge 0}$ be a standard Brownian motion and let
\[
M_t:=\sup_{0\le s\le t} B_s.
\]
Then, for every $a>0$ and every $t>0$,
\[
\mathbb{P}(M_t\ge a)
=
2\,\mathbb{P}(B_t\ge a).
\]
Equivalently,
\[
\mathbb{P}\left(\sup_{0\le s\le t}B_s\ge a\right)
=
2\left(1-\Phi\left(\frac{a}{\sqrt t}\right)\right),
\]
where $\Phi$ is the standard normal cumulative distribution function.
\end{theorem}


Let
\[
\tau_b^0:=\inf\{t\ge0:\widetilde B_t=b\},
\qquad b>0.
\]
We recall the standard density of the first hitting time of level $b$ by a Brownian motion:
\[
f_{\tau_b^0}(t)
=
\frac{b}{\sqrt{2\pi t^3}}
\exp\left(-\frac{b^2}{2t}\right),
\qquad t>0.
\]

Let us briefly derive it from the reflection principle. For $t>0$,
\[
\{\tau_b^0\le t\}
=
\left\{\sup_{0\le s\le t}\widetilde B_s\ge b\right\}.
\]
By the reflection principle,
\[
\mathbb{Q}(\tau_b^0\le t)
=
\mathbb{Q}\left(\sup_{0\le s\le t}\widetilde B_s\ge b\right)
=
2\mathbb{Q}(\widetilde B_t\ge b).
\]
Since $\widetilde B_t\sim \mathcal{N}(0,t)$ under $\mathbb{Q}$,
\[
\mathbb{Q}(\tau_b^0\le t)
=
2\left(1-\Phi\left(\frac{b}{\sqrt t}\right)\right).
\]
Differentiating with respect to $t$ gives the density:
\[
f_{\tau_b^0}(t)
=
2\frac{\dd}{\dd t}
\left[
1-\Phi\left(\frac{b}{\sqrt t}\right)
\right].
\]
Now
\[
\frac{\dd}{\dd t}\left(\frac{b}{\sqrt t}\right)
=
-\frac{b}{2t^{3/2}},
\]
and
\[
\Phi'(x)=\frac{1}{\sqrt{2\pi}}e^{-x^2/2}.
\]
Thus,
\[
f_{\tau_b^0}(t)
=
2
\left[
-\Phi'\left(\frac{b}{\sqrt t}\right)
\left(-\frac{b}{2t^{3/2}}\right)
\right]
=
\frac{b}{\sqrt{2\pi t^3}}
\exp\left(-\frac{b^2}{2t}\right).
\]
\end{solutionpart}

\begin{solutionpart}{Question (e)}
We now compute the density of $\tau_a$ under the original measure $\mathbb{P}$.

From Question (c), under $\mathbb{Q}$, $\tau_a$ has the same law as the first hitting time of level
\[
b=\frac{a}{\sigma}
\]
by a standard Brownian motion. Hence its density under $\mathbb{Q}$ is
\[
f_{\tau_a}^{\mathbb{Q}}(t)
=
\frac{a/\sigma}{\sqrt{2\pi t^3}}
\exp\left(-\frac{a^2}{2\sigma^2t}\right),
\qquad t>0.
\]

We now transform back to $\mathbb{P}$. Since
\[
\frac{\dd\mathbb{Q}}{\dd\mathbb{P}}\Big|_{\mathcal{F}_t}
=
Z_t,
\]
we have
\[
\frac{\dd\mathbb{P}}{\dd\mathbb{Q}}\Big|_{\mathcal{F}_t}
=
Z_t^{-1}.
\]
Moreover,
\[
Z_t^{-1}
=
\exp\left(\theta B_t+\frac12\theta^2t\right).
\]

On the event $\{\tau_a=t\}$, one has
\[
X_t=a.
\]
Since
\[
X_t=\mu t+\sigma B_t,
\]
this gives
\[
B_t=\frac{a-\mu t}{\sigma}.
\]
Therefore,
\[
Z_t^{-1}
=
\exp\left(
\theta\frac{a-\mu t}{\sigma}
+
\frac12\theta^2t
\right).
\]
Using $\theta=\mu/\sigma$, we get
\[
Z_t^{-1}
=
\exp\left(
\frac{\mu(a-\mu t)}{\sigma^2}
+
\frac12\frac{\mu^2}{\sigma^2}t
\right).
\]
Hence
\[
Z_t^{-1}
=
\exp\left(
\frac{\mu a}{\sigma^2}
-
\frac12\frac{\mu^2}{\sigma^2}t
\right).
\]

Thus, informally but rigorously justified by applying the change of measure to the event
$\{\tau_a\in \dd t\}$, we obtain
\[
f_{\tau_a}^{\mathbb{P}}(t)
=
\exp\left(
\frac{\mu a}{\sigma^2}
-
\frac12\frac{\mu^2}{\sigma^2}t
\right)
f_{\tau_a}^{\mathbb{Q}}(t).
\]
Therefore,
\[
f_{\tau_a}^{\mathbb{P}}(t)
=
\exp\left(
\frac{\mu a}{\sigma^2}
-
\frac12\frac{\mu^2}{\sigma^2}t
\right)
\frac{a/\sigma}{\sqrt{2\pi t^3}}
\exp\left(-\frac{a^2}{2\sigma^2t}\right).
\]
\end{solutionpart}

\begin{solutionpart}{Question (f)}
We simplify the expression obtained above:
\[
f_{\tau_a}^{\mathbb{P}}(t)
=
\frac{a}{\sigma\sqrt{2\pi t^3}}
\exp\left(
-\frac{a^2}{2\sigma^2t}
+
\frac{\mu a}{\sigma^2}
-
\frac{\mu^2 t}{2\sigma^2}
\right).
\]
We now combine the terms in the exponential. Observe that
\[
-\frac{(a-\mu t)^2}{2\sigma^2 t}
=
-\frac{a^2-2\mu at+\mu^2t^2}{2\sigma^2t}
=
-\frac{a^2}{2\sigma^2t}
+
\frac{\mu a}{\sigma^2}
-
\frac{\mu^2t}{2\sigma^2}.
\]
Therefore,
\[
f_{\tau_a}^{\mathbb{P}}(t)
=
\frac{a}{\sigma\sqrt{2\pi t^3}}
\exp\left(
-\frac{(a-\mu t)^2}{2\sigma^2t}
\right),
\qquad t>0.
\]

Thus,
\[
\boxed{
f_{\tau_a}(t)
=
\frac{a}{\sigma\sqrt{2\pi t^3}}
\exp\left(
-\frac{(a-\mu t)^2}{2\sigma^2t}
\right),
\qquad t>0.
}
\]
This is the density of an inverse Gaussian distribution. More precisely,
\[
\tau_a
\sim
\operatorname{IG}\left(
\frac{a}{\mu},
\frac{a^2}{\sigma^2}
\right),
\]
with the convention that an inverse Gaussian random variable with parameters
$(m,\lambda)$ has density
\[
f(t)
=
\left(\frac{\lambda}{2\pi t^3}\right)^{1/2}
\exp\left(
-\frac{\lambda(t-m)^2}{2m^2t}
\right),
\qquad t>0.
\]
\end{solutionpart}
\end{solution}

% ============================================================
\section{Martingale Pricing}
% ============================================================

% To be completed.
\begin{solution}
\textbf{Asian option: arithmetic average}
We work under the risk-neutral measure $\mathbb{Q}$. The stock price follows
\[
\frac{\dd S_t}{S_t}=r\,\dd t+\sigma\,\dd B_t,
\qquad S_0>0,
\]
or equivalently
\[
\dd S_t=rS_t\,\dd t+\sigma S_t\,\dd B_t.
\]
The payoff of the arithmetic-average Asian call is
\[
H=
\left(
\frac1T\int_0^T S_u\,\dd u-K
\right)^+.
\]
We define the running integral
\[
I_t:=\int_0^t S_u\,\dd u.
\]

\begin{solutionpart}{Question (a)}
By the risk-neutral pricing formula, the value at time $t$ is the discounted conditional expectation of the payoff:
\[
V_t
=
\E^\mathbb{Q}\left[
e^{-r(T-t)}
\left(
\frac1T\int_0^T S_u\,\dd u-K
\right)^+
\Bigm|\F_t
\right].
\]
Since $r$ is constant, the discount factor is deterministic and can also be written outside the conditional expectation:
\[
V_t
=
e^{-r(T-t)}
\E^\mathbb{Q}\left[
\left(
\frac1T\int_0^T S_u\,\dd u-K
\right)^+
\Bigm|\F_t
\right].
\]
\end{solutionpart}

\begin{solutionpart}{Question (b)}
For any $t\in[0,T]$, we decompose the integral as
\[
\int_0^T S_u\,\dd u
=
\int_0^t S_u\,\dd u
+
\int_t^T S_u\,\dd u.
\]
By definition of $I_t$, this gives
\[
\int_0^T S_u\,\dd u
=
I_t+\int_t^T S_u\,\dd u.
\]

The future evolution after time $t$ depends on the past only through the current stock value $S_t$, because $S$ is a Markov diffusion. However, the payoff also depends on the accumulated past average, summarized by $I_t$.

Therefore, the price can be written as a deterministic function of the two state variables $(S_t,I_t)$:
\[
V_t=v(t,S_t,I_t),
\]
where
\[
v(t,s,i)
:=
\E_{t,s,i}^{\mathbb{Q}}\left[
e^{-r(T-t)}
\left(
\frac1T
\left(
i+\int_t^T S_u\,\dd u
\right)
-K
\right)^+
\right].
\]
Here $\E_{t,s,i}^{\mathbb{Q}}$ means that the process starts at time $t$ from
\[
S_t=s,
\qquad
I_t=i.
\]
\end{solutionpart}

\begin{solutionpart}{Question (c)}
The stock dynamics under $\mathbb{Q}$ are
\[
\dd S_t=rS_t\,\dd t+\sigma S_t\,\dd B_t.
\]
Since
\[
I_t=\int_0^t S_u\,\dd u,
\]
we have
\[
\dd I_t=S_t\,\dd t.
\]
Therefore, the two-dimensional process $(S_t,I_t)$ satisfies
\[
\boxed{
\begin{aligned}
\dd S_t &= rS_t\,\dd t+\sigma S_t\,\dd B_t,\\
\dd I_t &= S_t\,\dd t.
\end{aligned}
}
\]

The process $(S_t,I_t)$ is Markov because its future dynamics depend only on the current values
$(S_t,I_t)$ and not on the full past trajectory. More precisely, given $(S_t,I_t)=(s,i)$, the future stock process is generated from $s$, and the future value of the running integral is obtained by adding
\[
\int_t^u S_\ell\,\dd \ell
\]
to the current value $i$.

By contrast, $S_t$ alone is not enough to price the Asian option, because the payoff depends on the historical average
\[
\frac1T\int_0^T S_u\,\dd u.
\]
At time $t$, two paths may have the same current stock price $S_t$ but different accumulated integrals $I_t$, leading to different option values.
\end{solutionpart}

\begin{solutionpart}{Question (d)}
Assume that
\[
V_t=v(t,S_t,I_t)
\]
for a sufficiently smooth function
\[
v\in C^{1,2,1}.
\]
Applying It\^o's formula to $v(t,S_t,I_t)$ gives
\[
\dd v(t,S_t,I_t)
=
\partial_t v\,\dd t
+
\partial_s v\,\dd S_t
+
\partial_i v\,\dd I_t
+
\frac12\partial_{ss}v\,\dd\langle S\rangle_t.
\]
There is no second-order term in $I$ because $I$ has finite variation, and no cross-variation term because
\[
\dd\langle S,I\rangle_t=0.
\]
Now
\[
\dd S_t=rS_t\,\dd t+\sigma S_t\,\dd B_t,
\qquad
\dd I_t=S_t\,\dd t,
\]
and
\[
\dd\langle S\rangle_t=\sigma^2S_t^2\,\dd t.
\]
Therefore,
\[
\begin{aligned}
\dd V_t
&=
\bigg(
\partial_t v
+
rS_t\partial_s v
+
S_t\partial_i v
+
\frac12\sigma^2S_t^2\partial_{ss}v
\bigg)\dd t
+
\sigma S_t\partial_s v\,\dd B_t.
\end{aligned}
\]

Under the risk-neutral measure, the discounted option price
\[
e^{-rt}V_t
\]
must be a martingale. Equivalently, the drift of $V_t$ must be equal to $rV_t\,\dd t$. Thus
\[
\partial_t v
+
rs\,\partial_s v
+
s\,\partial_i v
+
\frac12\sigma^2s^2\partial_{ss}v
=
rv.
\]
Hence the PDE is
\[
\boxed{
\partial_t v
+
rs\,\partial_s v
+
s\,\partial_i v
+
\frac12\sigma^2s^2\partial_{ss}v
-
rv
=0.
}
\]

At maturity,
\[
V_T=H
=
\left(
\frac1T\int_0^T S_u\,\dd u-K
\right)^+
=
\left(
\frac{I_T}{T}-K
\right)^+.
\]
Thus the terminal condition is
\[
\boxed{
v(T,s,i)=\left(\frac{i}{T}-K\right)^+.
}
\]
\end{solutionpart}

\begin{solutionpart}{Question (e)}
Define the discounted price process
\[
\bar V_t:=e^{-rt}V_t.
\]
Using the risk-neutral pricing formula,
\[
V_t
=
\E^\mathbb{Q}\left[
e^{-r(T-t)}H
\Bigm|\F_t
\right].
\]
Multiplying both sides by $e^{-rt}$ gives
\[
e^{-rt}V_t
=
e^{-rt}
\E^\mathbb{Q}\left[
e^{-r(T-t)}H
\Bigm|\F_t
\right].
\]
Since
\[
e^{-rt}e^{-r(T-t)}=e^{-rT},
\]
we get
\[
\bar V_t
=
\E^\mathbb{Q}\left[
e^{-rT}H
\Bigm|\F_t
\right].
\]
Therefore,
\[
\bar V_t
=
\E^\mathbb{Q}\left[
e^{-rT}
\left(
\frac1T\int_0^T S_u\,\dd u-K
\right)^+
\Bigm|\F_t
\right].
\]
This is a conditional expectation of a fixed integrable random variable. Hence, by the tower property, $(\bar V_t)_{0\le t\le T}$ is a $\mathbb{Q}$-martingale.

Indeed, for $0\le s\le t\le T$,
\[
\E^\mathbb{Q}\left[\bar V_t\mid \F_s\right]
=
\E^\mathbb{Q}\left[
\E^\mathbb{Q}\left[
e^{-rT}H
\mid \F_t
\right]
\Bigm|\F_s
\right]
=
\E^\mathbb{Q}\left[
e^{-rT}H
\mid \F_s
\right]
=
\bar V_s.
\]
\end{solutionpart}

\begin{solutionpart}{Question (f)}
We propose a Monte Carlo estimator for $V_0$.

Let
\[
0=t_0<t_1<\cdots<t_n=T,
\qquad
\Delta t=\frac{T}{n}.
\]
Under $\mathbb{Q}$, simulate $M$ independent paths of the Black--Scholes dynamics
\[
\dd S_t=rS_t\,\dd t+\sigma S_t\,\dd B_t.
\]
For example, using the exact discretization,
\[
S_{t_{k+1}}^{(m)}
=
S_{t_k}^{(m)}
\exp\left(
\left(r-\frac12\sigma^2\right)\Delta t
+
\sigma\sqrt{\Delta t}\,Z_{k+1}^{(m)}
\right),
\]
where
\[
Z_{k+1}^{(m)}\sim \mathcal{N}(0,1)
\]
are independent.

For each simulated path, approximate the time integral by a Riemann sum:
\[
\int_0^T S_u^{(m)}\,\dd u
\approx
\sum_{k=0}^{n-1} S_{t_k}^{(m)}\Delta t.
\]
The corresponding simulated payoff is
\[
H^{(m)}_n
:=
\left(
\frac1T
\sum_{k=0}^{n-1} S_{t_k}^{(m)}\Delta t
-K
\right)^+.
\]
Then the Monte Carlo estimator is
\[
\boxed{
\widehat V_0^{\,M,n}
=
e^{-rT}
\frac1M\sum_{m=1}^M H_n^{(m)}.
}
\]

Thus:
\begin{itemize}[leftmargin=1.4em]
\item We simulate the stock price process under the risk-neutral measure $\mathbb{Q}$.
\item We average the simulated discounted payoffs.
\item The discount factor $e^{-rT}$ appears outside the Monte Carlo average, since $r$ is constant.
\end{itemize}

If one uses the Euler scheme instead of the exact Black--Scholes discretization, one simulates
\[
S_{t_{k+1}}^{(m)}
=
S_{t_k}^{(m)}
+
rS_{t_k}^{(m)}\Delta t
+
\sigma S_{t_k}^{(m)}\sqrt{\Delta t}\,Z_{k+1}^{(m)}.
\]
However, for a geometric Brownian motion, the exact lognormal discretization is usually preferred because it preserves positivity.
\end{solutionpart}

\begin{solutionpart}{Question (g)}
The geometric average is defined by
\[
G
:=
\exp\left(
\frac1T\int_0^T \log S_u\,\dd u
\right).
\]
In the Black--Scholes model,
\[
S_t
=
S_0
\exp\left(
\left(r-\frac12\sigma^2\right)t+\sigma B_t
\right).
\]
Therefore,
\[
\log S_t
=
\log S_0
+
\left(r-\frac12\sigma^2\right)t
+
\sigma B_t.
\]
Hence
\[
\frac1T\int_0^T \log S_u\,\dd u
=
\log S_0
+
\left(r-\frac12\sigma^2\right)\frac{1}{T}\int_0^T u\,\dd u
+
\frac{\sigma}{T}\int_0^T B_u\,\dd u.
\]
Since
\[
\int_0^T u\,\dd u=\frac{T^2}{2},
\]
we get
\[
\frac1T\int_0^T \log S_u\,\dd u
=
\log S_0
+
\frac12\left(r-\frac12\sigma^2\right)T
+
\frac{\sigma}{T}\int_0^T B_u\,\dd u.
\]
The random variable
\[
\int_0^T B_u\,\dd u
\]
is Gaussian because it is a linear functional of the Brownian path. Its mean is $0$ and its variance is
\[
\Var\left(\int_0^T B_u\,\dd u\right)
=
\int_0^T\int_0^T \Cov(B_u,B_v)\,\dd u\,\dd v
=
\int_0^T\int_0^T \min(u,v)\,\dd u\,\dd v.
\]
Now
\[
\int_0^T\int_0^T \min(u,v)\,\dd u\,\dd v
=
2\int_0^T\int_0^v u\,\dd u\,\dd v
=
2\int_0^T \frac{v^2}{2}\,\dd v
=
\frac{T^3}{3}.
\]
Therefore,
\[
\frac{\sigma}{T}\int_0^T B_u\,\dd u
\sim
\mathcal{N}\left(0,\frac{\sigma^2T}{3}\right).
\]

Thus,
\[
\log G
\sim
\mathcal{N}\left(
\log S_0+\frac12\left(r-\frac12\sigma^2\right)T,
\frac{\sigma^2T}{3}
\right).
\]
Hence $G$ is lognormally distributed. This is why the geometric-average Asian option is analytically tractable in the Black--Scholes model.

In particular, a call option with payoff
\[
(G-K)^+
\]
can be priced using a Black--Scholes type formula with an adjusted mean and volatility.

This can be used as a control variate for the arithmetic Asian option. Let
\[
H_A
=
\left(
\frac1T\int_0^T S_u\,\dd u-K
\right)^+
\]
be the arithmetic Asian payoff and
\[
H_G=(G-K)^+
\]
be the geometric Asian payoff. Since the price of $H_G$ is known explicitly, one can estimate
\[
V_A
=
e^{-rT}\E^\mathbb{Q}[H_A]
\]
using
\[
\widehat V_A^{\,CV}
=
e^{-rT}
\left[
\frac1M\sum_{m=1}^M
\left(
H_A^{(m)}
-
\beta\left(H_G^{(m)}-\E^\mathbb{Q}[H_G]\right)
\right)
\right],
\]
where $\beta$ is chosen to reduce variance, ideally
\[
\beta^\star
=
\frac{\Cov(H_A,H_G)}{\Var(H_G)}.
\]
Because the arithmetic and geometric averages are highly correlated, this control variate is usually very effective.
\end{solutionpart}
\end{solution}

% ------------------------------------------------------------


\begin{solution}
\textbf{Continuous-time lookback digital one-touch in Black--Scholes}
We work under the risk-neutral measure $\mathbb{Q}$. The stock price satisfies
\[
\frac{\dd S_t}{S_t}=(r-q)\,\dd t+\sigma\,\dd B_t,
\qquad S_0>0,\quad \sigma>0.
\]
The payoff is
\[
\Pi=\mathbf 1_{\left\{\sup_{0\le t\le T}S_t\ge H\right\}},
\qquad H>S_0.
\]
Hence the no-arbitrage price is
\[
V_0=e^{-rT}\mathbb{Q}\left(\sup_{0\le t\le T}S_t\ge H\right).
\]

\begin{solutionpart}{Question (a)}
Define
\[
X_t:=\ln\left(\frac{S_t}{S_0}\right).
\]
Since
\[
\dd S_t=(r-q)S_t\,\dd t+\sigma S_t\,\dd B_t,
\]
It\^o's formula applied to $\log S_t$ gives
\[
\dd \log S_t
=
\frac{1}{S_t}\,\dd S_t
-
\frac12\frac{1}{S_t^2}\,\dd\langle S\rangle_t.
\]
Now
\[
\dd\langle S\rangle_t=\sigma^2 S_t^2\,\dd t.
\]
Therefore,
\[
\dd \log S_t
=
(r-q)\,\dd t+\sigma\,\dd B_t
-
\frac12\sigma^2\,\dd t.
\]
Thus
\[
\dd \log S_t
=
\left(r-q-\frac12\sigma^2\right)\dd t+\sigma\,\dd B_t.
\]
Integrating from $0$ to $t$ gives
\[
\log S_t-\log S_0
=
\left(r-q-\frac12\sigma^2\right)t+\sigma B_t.
\]
Hence
\[
X_t
=
\ln\left(\frac{S_t}{S_0}\right)
=
\left(r-q-\frac12\sigma^2\right)t+\sigma B_t.
\]
If we define
\[
\mu:=r-q-\frac12\sigma^2,
\]
then
\[
\boxed{
X_t=\mu t+\sigma B_t.
}
\]

Now, since the logarithm is increasing,
\[
S_t\ge H
\quad\Longleftrightarrow\quad
\ln\left(\frac{S_t}{S_0}\right)\ge \ln\left(\frac{H}{S_0}\right).
\]
Define
\[
a:=\ln\left(\frac{H}{S_0}\right).
\]
Since $H>S_0$, we have $a>0$. Therefore,
\[
\left\{\sup_{0\le t\le T}S_t\ge H\right\}
=
\left\{\sup_{0\le t\le T}X_t\ge a\right\}.
\]
So the pricing problem is reduced to computing
\[
\mathbb{Q}\left(\sup_{0\le t\le T}(\mu t+\sigma B_t)\ge a\right).
\]
\end{solutionpart}

\begin{solutionpart}{Question (b)}
We now prove the hitting probability formula for the drifted Brownian motion
\[
Y_t:=\mu t+\sigma B_t.
\]
We want to show that, for $a>0$,
\[
\mathbb{Q}\left(\sup_{0\le t\le T}Y_t\ge a\right)
=
\Phi\left(\frac{\mu T-a}{\sigma\sqrt T}\right)
+
\exp\left(\frac{2\mu a}{\sigma^2}\right)
\Phi\left(\frac{-\mu T-a}{\sigma\sqrt T}\right).
\]

\medskip

First, consider the driftless case $\mu=0$. Then
\[
Y_t=\sigma B_t.
\]
By the reflection principle,
\[
\mathbb{Q}\left(\sup_{0\le t\le T}\sigma B_t\ge a\right)
=
2\mathbb{Q}(\sigma B_T\ge a).
\]
Since $B_T\sim \mathcal N(0,T)$,
\[
\sigma B_T\sim \mathcal N(0,\sigma^2T),
\]
and therefore
\[
\mathbb{Q}(\sigma B_T\ge a)
=
1-\Phi\left(\frac{a}{\sigma\sqrt T}\right).
\]
Thus
\[
\mathbb{Q}\left(\sup_{0\le t\le T}\sigma B_t\ge a\right)
=
2\left[
1-\Phi\left(\frac{a}{\sigma\sqrt T}\right)
\right].
\]
Using the identity $1-\Phi(x)=\Phi(-x)$, this becomes
\[
\mathbb{Q}\left(\sup_{0\le t\le T}\sigma B_t\ge a\right)
=
2\Phi\left(-\frac{a}{\sigma\sqrt T}\right).
\]
This is exactly the previous formula when $\mu=0$, since
\[
\Phi\left(\frac{-a}{\sigma\sqrt T}\right)
+
\Phi\left(\frac{-a}{\sigma\sqrt T}\right)
=
2\Phi\left(-\frac{a}{\sigma\sqrt T}\right).
\]

\medskip

Now assume $\mu\neq 0$. Define
\[
\theta:=\frac{\mu}{\sigma}
\]
and
\[
Z_T:=
\exp\left(
-\theta B_T-\frac12\theta^2T
\right).
\]
Define a new probability measure $\widetilde{\mathbb{Q}}$ by
\[
\frac{\dd\widetilde{\mathbb{Q}}}{\dd\mathbb{Q}}
=
Z_T.
\]
By Girsanov's theorem,
\[
\widetilde B_t:=B_t+\theta t
\]
is a Brownian motion under $\widetilde{\mathbb{Q}}$.

Since
\[
B_t=\widetilde B_t-\theta t,
\]
we have
\[
Y_t
=
\mu t+\sigma B_t
=
\mu t+\sigma(\widetilde B_t-\theta t).
\]
Using $\theta=\mu/\sigma$, we obtain
\[
Y_t=\sigma\widetilde B_t.
\]
Thus, under $\widetilde{\mathbb{Q}}$, the event
\[
\left\{\sup_{0\le t\le T}Y_t\ge a\right\}
\]
becomes
\[
\left\{\sup_{0\le t\le T}\sigma\widetilde B_t\ge a\right\}.
\]

However, we need the probability of this event under the original measure $\mathbb{Q}$.
Let
\[
\tau_a:=\inf\{t\ge0:Y_t=a\}.
\]
Then
\[
\left\{\sup_{0\le t\le T}Y_t\ge a\right\}
=
\{\tau_a\le T\}.
\]
Using the change of measure,
\[
\mathbb{Q}(\tau_a\le T)
=
\mathbb{E}^{\widetilde{\mathbb{Q}}}
\left[
Z_T^{-1}\mathbf 1_{\{\tau_a\le T\}}
\right].
\]
Now
\[
Z_T^{-1}
=
\exp\left(
\theta B_T+\frac12\theta^2T
\right).
\]
Since
\[
B_T=\widetilde B_T-\theta T,
\]
we get
\[
Z_T^{-1}
=
\exp\left(
\theta\widetilde B_T-\frac12\theta^2T
\right).
\]
Therefore,
\[
\mathbb{Q}(\tau_a\le T)
=
\mathbb{E}^{\widetilde{\mathbb{Q}}}
\left[
\exp\left(
\theta\widetilde B_T-\frac12\theta^2T
\right)
\mathbf 1_{\{\sup_{0\le t\le T}\sigma\widetilde B_t\ge a\}}
\right].
\]

Let
\[
b:=\frac{a}{\sigma}.
\]
Then
\[
\{\sup_{0\le t\le T}\sigma\widetilde B_t\ge a\}
=
\{\sup_{0\le t\le T}\widetilde B_t\ge b\}.
\]
Hence
\[
\mathbb{Q}(\tau_a\le T)
=
e^{-\frac12\theta^2T}
\mathbb{E}^{\widetilde{\mathbb{Q}}}
\left[
e^{\theta\widetilde B_T}
\mathbf 1_{\{\sup_{0\le t\le T}\widetilde B_t\ge b\}}
\right].
\]

We now use the reflection principle in its density form. For Brownian motion,
on the event that the maximum crosses $b$, the endpoint distribution is reflected across $b$.
Thus
\[
\mathbb{E}^{\widetilde{\mathbb{Q}}}
\left[
e^{\theta\widetilde B_T}
\mathbf 1_{\{\sup_{0\le t\le T}\widetilde B_t\ge b\}}
\right]
=
\int_b^\infty e^{\theta x}\frac{1}{\sqrt{2\pi T}}e^{-\frac{x^2}{2T}}\,\dd x
+
\int_{-\infty}^b e^{\theta x}
\frac{1}{\sqrt{2\pi T}}
e^{-\frac{(2b-x)^2}{2T}}\,\dd x.
\]
Call these two terms $I_1$ and $I_2$.

For the first integral,
\[
I_1
=
\int_b^\infty e^{\theta x}\frac{1}{\sqrt{2\pi T}}e^{-\frac{x^2}{2T}}\,\dd x.
\]
Complete the square:
\[
\theta x-\frac{x^2}{2T}
=
-\frac{(x-\theta T)^2}{2T}
+
\frac12\theta^2T.
\]
Hence
\[
I_1
=
e^{\frac12\theta^2T}
\int_b^\infty
\frac{1}{\sqrt{2\pi T}}
e^{-\frac{(x-\theta T)^2}{2T}}\,\dd x.
\]
Therefore
\[
I_1
=
e^{\frac12\theta^2T}
\Phi\left(\frac{\theta T-b}{\sqrt T}\right).
\]

For the second integral,
\[
I_2
=
\int_{-\infty}^b e^{\theta x}
\frac{1}{\sqrt{2\pi T}}
e^{-\frac{(2b-x)^2}{2T}}\,\dd x.
\]
Make the change of variables
\[
y=2b-x,
\qquad
\dd y=-\dd x.
\]
When $x\le b$, we have $y\ge b$. Hence
\[
I_2
=
\int_b^\infty
e^{\theta(2b-y)}
\frac{1}{\sqrt{2\pi T}}
e^{-\frac{y^2}{2T}}\,\dd y.
\]
Thus
\[
I_2
=
e^{2\theta b}
\int_b^\infty
e^{-\theta y}
\frac{1}{\sqrt{2\pi T}}
e^{-\frac{y^2}{2T}}\,\dd y.
\]
Again complete the square:
\[
-\theta y-\frac{y^2}{2T}
=
-\frac{(y+\theta T)^2}{2T}
+
\frac12\theta^2T.
\]
Therefore
\[
I_2
=
e^{2\theta b}
e^{\frac12\theta^2T}
\int_b^\infty
\frac{1}{\sqrt{2\pi T}}
e^{-\frac{(y+\theta T)^2}{2T}}\,\dd y.
\]
This gives
\[
I_2
=
e^{2\theta b}
e^{\frac12\theta^2T}
\Phi\left(
\frac{-b-\theta T}{\sqrt T}
\right).
\]

Combining $I_1$ and $I_2$, we find
\[
\begin{aligned}
\mathbb{Q}(\tau_a\le T)
&=
e^{-\frac12\theta^2T}(I_1+I_2)\\
&=
\Phi\left(\frac{\theta T-b}{\sqrt T}\right)
+
e^{2\theta b}
\Phi\left(\frac{-b-\theta T}{\sqrt T}\right).
\end{aligned}
\]
Recalling that
\[
\theta=\frac{\mu}{\sigma},
\qquad
b=\frac{a}{\sigma},
\]
we get
\[
\frac{\theta T-b}{\sqrt T}
=
\frac{\mu T-a}{\sigma\sqrt T},
\]
and
\[
\frac{-b-\theta T}{\sqrt T}
=
\frac{-a-\mu T}{\sigma\sqrt T}.
\]
Moreover,
\[
2\theta b
=
2\frac{\mu}{\sigma}\frac{a}{\sigma}
=
\frac{2\mu a}{\sigma^2}.
\]
Therefore,
\[
\boxed{
\mathbb{Q}\left(\sup_{0\le t\le T}Y_t\ge a\right)
=
\Phi\left(\frac{\mu T-a}{\sigma\sqrt T}\right)
+
\exp\left(\frac{2\mu a}{\sigma^2}\right)
\Phi\left(\frac{-\mu T-a}{\sigma\sqrt T}\right).
}
\]
\end{solutionpart}

\begin{solutionpart}{Question (c)}
From Question (a), we have
\[
\left\{\sup_{0\le t\le T}S_t\ge H\right\}
=
\left\{\sup_{0\le t\le T}X_t\ge a\right\},
\]
where
\[
X_t=\mu t+\sigma B_t,
\qquad
\mu=r-q-\frac12\sigma^2,
\qquad
a=\ln\left(\frac{H}{S_0}\right).
\]
Using the hitting probability formula from Question (b), we obtain
\[
\mathbb{Q}\left(\sup_{0\le t\le T}S_t\ge H\right)
=
\Phi\left(\frac{\mu T-a}{\sigma\sqrt T}\right)
+
\exp\left(\frac{2\mu a}{\sigma^2}\right)
\Phi\left(\frac{-\mu T-a}{\sigma\sqrt T}\right).
\]
Therefore the time-$0$ price is
\[
\boxed{
V_0
=
e^{-rT}
\left[
\Phi\left(\frac{\mu T-a}{\sigma\sqrt T}\right)
+
\exp\left(\frac{2\mu a}{\sigma^2}\right)
\Phi\left(\frac{-\mu T-a}{\sigma\sqrt T}\right)
\right],
}
\]
where
\[
\mu=r-q-\frac12\sigma^2,
\qquad
a=\ln\left(\frac{H}{S_0}\right).
\]

Let us check the limiting cases.

If $H\downarrow S_0$, then
\[
a=\ln(H/S_0)\downarrow0.
\]
The event
\[
\left\{\sup_{0\le t\le T}S_t\ge S_0\right\}
\]
has probability $1$, since $S_0$ is already reached at time $0$. Thus
\[
V_0\to e^{-rT}.
\]
This is also consistent with the formula, since when $a=0$,
\[
\Phi\left(\frac{\mu T}{\sigma\sqrt T}\right)
+
\Phi\left(\frac{-\mu T}{\sigma\sqrt T}\right)
=
\Phi(x)+\Phi(-x)=1.
\]

If $H\uparrow\infty$, then
\[
a=\ln(H/S_0)\to\infty.
\]
The probability of hitting such a large barrier before the fixed maturity $T$ tends to $0$.
Therefore,
\[
V_0\to0.
\]
\end{solutionpart}

\begin{solutionpart}{Question (d)}
For a Monte Carlo verification, simulate the stock price under $\mathbb{Q}$ on a grid
\[
0=t_0<t_1<\cdots<t_N=T.
\]
Using the exact Black--Scholes discretization,
\[
S_{t_{k+1}}^{(m)}
=
S_{t_k}^{(m)}
\exp\left(
\left(r-q-\frac12\sigma^2\right)\Delta t
+
\sigma\sqrt{\Delta t}\,Z_{k+1}^{(m)}
\right),
\]
where
\[
Z_{k+1}^{(m)}\sim\mathcal N(0,1)
\]
are independent.

For each path, approximate the event
\[
\left\{\sup_{0\le t\le T}S_t\ge H\right\}
\]
by the discretely monitored event
\[
\left\{\max_{0\le k\le N}S_{t_k}^{(m)}\ge H\right\}.
\]
The Monte Carlo estimator is
\[
\boxed{
\widehat V_0^{(N,M)}
=
e^{-rT}\frac1M
\sum_{m=1}^M
\mathbf 1_{\left\{\max_{0\le k\le N}S_{t_k}^{(m)}\ge H\right\}}.
}
\]

This estimator converges, as $M\to\infty$, to the price of the discretely monitored one-touch option. It is generally biased downward relative to the continuously monitored price, because the simulated grid may miss barrier crossings occurring between two consecutive time points.

As $N$ increases, the time grid becomes finer, and the discretely monitored event better approximates the continuously monitored event. Hence the bias decreases as $N\to\infty$.
\end{solutionpart}
\end{solution}
% ============================================================
\section{American Options}
% ============================================================

\begin{solution}
We consider a three-period binomial model with
\[
S_0=100,
\qquad
u=1.1,
\qquad
d=0.9,
\qquad
1+r=1.02,
\]
and strike
\[
K=100.
\]

The risk-neutral probability $p$ is determined by
\[
S_0
=
\frac{1}{1+r}
\left(
p\,uS_0+(1-p)dS_0
\right).
\]
Equivalently,
\[
1+r=pu+(1-p)d.
\]
Thus
\[
1.02=1.1p+0.9(1-p).
\]
Hence
\[
1.02=0.9+0.2p,
\]
so
\[
p=\frac{1.02-0.9}{1.1-0.9}=0.6.
\]
Therefore,
\[
\boxed{p=0.6,\qquad 1-p=0.4.}
\]

\begin{solutionpart}{Question (a): European put price}

At maturity $T=3$, the possible stock prices are
\[
S_3^{uuu}=100(1.1)^3=133.1,
\]
\[
S_3^{uud}=100(1.1)^2(0.9)=108.9,
\]
\[
S_3^{udd}=100(1.1)(0.9)^2=89.1,
\]
and
\[
S_3^{ddd}=100(0.9)^3=72.9.
\]

The European put payoff is
\[
(K-S_3)^+.
\]
Therefore, the terminal payoffs are
\[
P_3^{uuu}=(100-133.1)^+=0,
\]
\[
P_3^{uud}=(100-108.9)^+=0,
\]
\[
P_3^{udd}=(100-89.1)^+=10.9,
\]
and
\[
P_3^{ddd}=(100-72.9)^+=27.1.
\]

The European put price is the discounted risk-neutral expectation:
\[
P_0^{\mathrm{Eur}}
=
\frac{1}{(1+r)^3}
\E^{\mathbb{Q}}\left[(K-S_3)^+\right].
\]
Using the binomial probabilities,
\[
\mathbb{Q}(uuu)=p^3,
\]
\[
\mathbb{Q}(uud)=3p^2(1-p),
\]
\[
\mathbb{Q}(udd)=3p(1-p)^2,
\]
and
\[
\mathbb{Q}(ddd)=(1-p)^3.
\]
Hence
\[
P_0^{\mathrm{Eur}}
=
\frac{1}{1.02^3}
\left[
p^3\cdot 0
+
3p^2(1-p)\cdot 0
+
3p(1-p)^2\cdot 10.9
+
(1-p)^3\cdot 27.1
\right].
\]
With $p=0.6$, this gives
\[
P_0^{\mathrm{Eur}}
=
\frac{1}{1.02^3}
\left[
3(0.6)(0.4)^2(10.9)
+
(0.4)^3(27.1)
\right].
\]
Now
\[
3(0.6)(0.4)^2(10.9)
=
3(0.6)(0.16)(10.9)
=
3.1392,
\]
and
\[
(0.4)^3(27.1)=0.064(27.1)=1.7344.
\]
Thus
\[
\E^{\mathbb{Q}}\left[(K-S_3)^+\right]
=
3.1392+1.7344
=
4.8736.
\]
Therefore,
\[
\boxed{
P_0^{\mathrm{Eur}}
=
\frac{4.8736}{1.02^3}
\approx 4.5918.
}
\]
\end{solutionpart}

\begin{solutionpart}{Question (b): American put price}

\begin{theorem}[Dynamic programming principle for American options]
Consider a discrete-time financial model with dates
\[
0=t_0<t_1<\cdots<t_N=T.
\]
Let $(\mathcal F_n)_{0\le n\le N}$ be the market filtration, and assume that under the risk-neutral measure
$\mathbb Q$, discounted asset prices are martingales.

Let $r$ be the constant interest rate per period, so that the one-period discount factor is
\[
\frac{1}{1+r}.
\]
Let the American payoff process be
\[
G_n:=g(n,S_n),
\qquad n=0,\dots,N.
\]
The value process $(V_n)_{0\le n\le N}$ of the American option is defined backward by
\[
V_N=G_N,
\]
and, for $n=N-1,\dots,0$,
\[
\boxed{
V_n
=
\max\left(
G_n,
\frac{1}{1+r}\,
\mathbb E^{\mathbb Q}\left[V_{n+1}\mid \mathcal F_n\right]
\right).
}
\]
Equivalently,
\[
\boxed{
V_n
=
\max\left(
\text{immediate exercise value},
\text{continuation value}
\right).
}
\]
\end{theorem}


For the American put, we compute the value by backward induction. At each node, the value is the maximum between:

\begin{itemize}[leftmargin=1.4em]
\item the immediate exercise value,
\[
(K-S)^+,
\]
\item the continuation value,
\[
\frac{1}{1+r}
\left(
pV^{u}+(1-p)V^{d}
\right).
\]
\end{itemize}

At maturity, the American and European values coincide:
\[
V_3^{uuu}=0,
\qquad
V_3^{uud}=0,
\qquad
V_3^{udd}=10.9,
\qquad
V_3^{ddd}=27.1.
\]

\medskip

\noindent\textbf{Step 2.}

At time $2$, the possible stock prices are
\[
S_2^{uu}=100(1.1)^2=121,
\]
\[
S_2^{ud}=100(1.1)(0.9)=99,
\]
and
\[
S_2^{dd}=100(0.9)^2=81.
\]

At node $uu$, the immediate exercise value is
\[
(100-121)^+=0.
\]
The continuation value is
\[
\frac{1}{1.02}
\left(
0.6\cdot 0+0.4\cdot 0
\right)=0.
\]
Therefore,
\[
V_2^{uu}=0.
\]

At node $ud$, the immediate exercise value is
\[
(100-99)^+=1.
\]
The continuation value is
\[
\frac{1}{1.02}
\left(
0.6\cdot 0+0.4\cdot 10.9
\right)
=
\frac{4.36}{1.02}
\approx 4.2745.
\]
Hence continuation is better, and
\[
V_2^{ud}
=
\max(1,4.2745)
=
4.2745.
\]

At node $dd$, the immediate exercise value is
\[
(100-81)^+=19.
\]
The continuation value is
\[
\frac{1}{1.02}
\left(
0.6\cdot 10.9+0.4\cdot 27.1
\right).
\]
Compute
\[
0.6\cdot 10.9+0.4\cdot 27.1
=
6.54+10.84
=
17.38.
\]
Thus the continuation value is
\[
\frac{17.38}{1.02}
\approx 17.0392.
\]
Since
\[
19>17.0392,
\]
immediate exercise is optimal. Therefore,
\[
V_2^{dd}=19.
\]

So at time $2$,
\[
\boxed{
V_2^{uu}=0,
\qquad
V_2^{ud}\approx 4.2745,
\qquad
V_2^{dd}=19.
}
\]

\medskip

\noindent\textbf{Step 1.}

At time $1$, the possible stock prices are
\[
S_1^u=100(1.1)=110,
\]
and
\[
S_1^d=100(0.9)=90.
\]

At node $u$, the immediate exercise value is
\[
(100-110)^+=0.
\]
The continuation value is
\[
\frac{1}{1.02}
\left(
0.6V_2^{uu}+0.4V_2^{ud}
\right)
=
\frac{1}{1.02}
\left(
0.6\cdot 0+0.4\cdot 4.2745
\right).
\]
Thus
\[
\text{continuation value}
=
\frac{1.7098}{1.02}
\approx 1.6763.
\]
Hence
\[
V_1^u
=
\max(0,1.6763)
=
1.6763.
\]

At node $d$, the immediate exercise value is
\[
(100-90)^+=10.
\]
The continuation value is
\[
\frac{1}{1.02}
\left(
0.6V_2^{ud}+0.4V_2^{dd}
\right).
\]
Therefore,
\[
\text{continuation value}
=
\frac{1}{1.02}
\left(
0.6\cdot 4.2745+0.4\cdot 19
\right).
\]
Compute
\[
0.6\cdot 4.2745+0.4\cdot 19
=
2.5647+7.6
=
10.1647.
\]
Thus
\[
\text{continuation value}
=
\frac{10.1647}{1.02}
\approx 9.9654.
\]
Since
\[
10>9.9654,
\]
immediate exercise is optimal. Therefore,
\[
V_1^d=10.
\]

So at time $1$,
\[
\boxed{
V_1^u\approx 1.6763,
\qquad
V_1^d=10.
}
\]

\medskip

\noindent\textbf{Step 0.}

At time $0$, the immediate exercise value is
\[
(100-100)^+=0.
\]
The continuation value is
\[
\frac{1}{1.02}
\left(
0.6V_1^u+0.4V_1^d
\right).
\]
Hence
\[
V_0^{\mathrm{Am}}
=
\frac{1}{1.02}
\left(
0.6\cdot 1.6763+0.4\cdot 10
\right).
\]
Compute
\[
0.6\cdot 1.6763+0.4\cdot 10
=
1.0058+4
=
5.0058.
\]
Therefore,
\[
\boxed{
V_0^{\mathrm{Am}}
=
\frac{5.0058}{1.02}
\approx 4.9077.
}
\]
\end{solutionpart}

\begin{solutionpart}{Question (c): Early exercise premium}

The European put price is
\[
P_0^{\mathrm{Eur}}\approx 4.5918.
\]
The American put price is
\[
P_0^{\mathrm{Am}}\approx 4.9077.
\]
Therefore, the early exercise premium is
\[
P_0^{\mathrm{Am}}-P_0^{\mathrm{Eur}}
\approx
4.9077-4.5918
=
0.3159.
\]
Thus
\[
\boxed{
P_0^{\mathrm{Am}}-P_0^{\mathrm{Eur}}
\approx 0.316.
}
\]

The American put is more valuable than the European put because the holder has the additional right to exercise before maturity. For a put option, early exercise can be optimal when the stock price is sufficiently low.
\end{solutionpart}

\begin{solutionpart}{Question (d): Optimal exercise nodes}

We identify the nodes where immediate exercise is optimal.

At time $2$:
\[
S_2^{uu}=121,
\qquad
S_2^{ud}=99,
\qquad
S_2^{dd}=81.
\]
We found:
\[
V_2^{uu}=0,
\]
with immediate exercise value $0$ and continuation value $0$.
This is an indifference node.

At node $ud$,
\[
S_2^{ud}=99,
\]
the immediate exercise value is $1$, while the continuation value is approximately $4.2745$. Therefore, continuation is optimal.

At node $dd$,
\[
S_2^{dd}=81,
\]
the immediate exercise value is $19$, while the continuation value is approximately $17.0392$. Therefore, immediate exercise is optimal.

At time $1$:
\[
S_1^u=110,
\qquad
S_1^d=90.
\]
At node $u$, continuation is optimal because the immediate exercise value is $0$ and the continuation value is approximately $1.6763$.

At node $d$,
\[
S_1^d=90,
\]
the immediate exercise value is $10$, while the continuation value is approximately $9.9654$. Therefore, immediate exercise is optimal.

At time $0$, the immediate exercise value is $0$, while the continuation value is approximately $4.9077$, so continuation is optimal.

Hence the strict early exercise nodes are
\[
\boxed{
S_1^d=90
\quad\text{and}\quad
S_2^{dd}=81.
}
\]
At the node
\[
S_2^{uu}=121,
\]
the holder is indifferent between exercise and continuation, since both values are $0$, but this is not economically a meaningful early exercise node.
\end{solutionpart}
\end{solution}
% ============================================================
\section{Numerics}
% ============================================================
\begin{solution}
This section recalls the basic numerical ideas used to simulate stochastic differential equations,
stochastic integrals, and Monte Carlo estimators.

\begin{solutionpart}{Euler--Maruyama scheme}

Consider a one-dimensional stochastic differential equation
\[
\dd X_t=b(t,X_t)\,\dd t+\sigma(t,X_t)\,\dd B_t,
\qquad X_0=x_0.
\]
Let
\[
0=t_0<t_1<\cdots<t_N=T,
\qquad
\Delta t=\frac{T}{N}.
\]
The Euler--Maruyama scheme is defined by
\[
\boxed{
X_{t_{k+1}}^{\Delta}
=
X_{t_k}^{\Delta}
+
b(t_k,X_{t_k}^{\Delta})\Delta t
+
\sigma(t_k,X_{t_k}^{\Delta})\Delta B_k,
}
\]
where
\[
\Delta B_k:=B_{t_{k+1}}-B_{t_k}.
\]
Since Brownian increments are independent and Gaussian, we simulate
\[
\Delta B_k=\sqrt{\Delta t}\,Z_k,
\qquad
Z_k\sim\mathcal N(0,1),
\]
with the variables $(Z_k)_{k=0}^{N-1}$ independent.

Thus the implementable form is
\[
\boxed{
X_{t_{k+1}}^{\Delta}
=
X_{t_k}^{\Delta}
+
b(t_k,X_{t_k}^{\Delta})\Delta t
+
\sigma(t_k,X_{t_k}^{\Delta})\sqrt{\Delta t}\,Z_k.
}
\]

For a geometric Brownian motion,
\[
\dd S_t=\mu S_t\,\dd t+\sigma S_t\,\dd B_t,
\]
Euler--Maruyama gives
\[
S_{t_{k+1}}^{\Delta}
=
S_{t_k}^{\Delta}
+
\mu S_{t_k}^{\Delta}\Delta t
+
\sigma S_{t_k}^{\Delta}\sqrt{\Delta t}\,Z_k.
\]
However, in the Black--Scholes model one often uses the exact discretization
\[
\boxed{
S_{t_{k+1}}
=
S_{t_k}
\exp\left(
\left(\mu-\frac12\sigma^2\right)\Delta t
+
\sigma\sqrt{\Delta t}\,Z_k
\right),
}
\]
which preserves positivity and has no time-discretization error for the marginal dynamics.
\end{solutionpart}

\begin{solutionpart}{Simulation of ordinary time integrals}

Many payoffs involve integrals of the form
\[
\int_0^T g(t,X_t)\,\dd t.
\]
Once a path has been simulated on the grid $(t_k)$, a simple left-point Riemann approximation is
\[
\boxed{
\int_0^T g(t,X_t)\,\dd t
\approx
\sum_{k=0}^{N-1} g(t_k,X_{t_k}^{\Delta})\Delta t.
}
\]
For example, for an arithmetic Asian option,
\[
\int_0^T S_t\,\dd t
\approx
\sum_{k=0}^{N-1} S_{t_k}^{\Delta}\Delta t.
\]
One could also use a trapezoidal approximation,
\[
\int_0^T g(t,X_t)\,\dd t
\approx
\sum_{k=0}^{N-1}
\frac{
g(t_k,X_{t_k}^{\Delta})
+
g(t_{k+1},X_{t_{k+1}}^{\Delta})
}{2}
\Delta t.
\]
\end{solutionpart}

\begin{solutionpart}{Simulation of It\^o stochastic integrals}

For an It\^o integral
\[
\int_0^T H_t\,\dd B_t,
\]
where $H_t$ is adapted, the natural discretization is the left-point approximation
\[
\boxed{
\int_0^T H_t\,\dd B_t
\approx
\sum_{k=0}^{N-1} H_{t_k}^{\Delta}\Delta B_k.
}
\]
Equivalently,
\[
\int_0^T H_t\,\dd B_t
\approx
\sum_{k=0}^{N-1}
H_{t_k}^{\Delta}\sqrt{\Delta t}\,Z_k.
\]

The left-point rule is important: it reflects the fact that the It\^o integral uses information available before the Brownian increment is realized. In particular,
\[
H_{t_k}^{\Delta}
\]
must be measurable with respect to the information available at time $t_k$.

For example, if
\[
H_t=f(t,X_t),
\]
then the It\^o integral is approximated by
\[
\boxed{
\int_0^T f(t,X_t)\,\dd B_t
\approx
\sum_{k=0}^{N-1}
f(t_k,X_{t_k}^{\Delta})\sqrt{\Delta t}\,Z_k.
}
\]
\end{solutionpart}

\begin{solutionpart}{Monte Carlo pricing}

Suppose a derivative has payoff $H$ at maturity $T$. Under the risk-neutral measure
$\mathbb Q$, its time-$0$ price is
\[
V_0=e^{-rT}\mathbb E^{\mathbb Q}[H].
\]
Monte Carlo consists in simulating $M$ independent paths under $\mathbb Q$, computing one payoff
for each path,
\[
H^{(1)},\dots,H^{(M)},
\]
and estimating the expectation by the empirical average:
\[
\boxed{
\widehat V_0
=
e^{-rT}\frac1M\sum_{m=1}^M H^{(m)}.
}
\]

For example, for an arithmetic Asian call,
\[
H
=
\left(
\frac1T\int_0^T S_t\,\dd t-K
\right)^+,
\]
we approximate
\[
H^{(m)}
\approx
\left(
\frac1T\sum_{k=0}^{N-1} S_{t_k}^{(m)}\Delta t-K
\right)^+,
\]
and use
\[
\boxed{
\widehat V_0
=
e^{-rT}
\frac1M\sum_{m=1}^M
\left(
\frac1T\sum_{k=0}^{N-1} S_{t_k}^{(m)}\Delta t-K
\right)^+.
}
\]

For a one-touch or lookback digital payoff,
\[
H=\mathbf 1_{\{\sup_{0\le t\le T}S_t\ge L\}},
\]
we approximate the supremum by the grid maximum:
\[
\sup_{0\le t\le T}S_t
\approx
\max_{0\le k\le N}S_{t_k}.
\]
The estimator becomes
\[
\boxed{
\widehat V_0
=
e^{-rT}
\frac1M\sum_{m=1}^M
\mathbf 1_{\left\{\max_{0\le k\le N}S_{t_k}^{(m)}\ge L\right\}}.
}
\]
This estimator has a discretization bias because the simulated grid can miss barrier crossings between two time points.
\end{solutionpart}

\begin{solutionpart}{Accuracy and sources of error}

There are two main sources of error.

First, there is a \emph{statistical Monte Carlo error}. Even if the payoff were simulated exactly,
\[
\frac1M\sum_{m=1}^M H^{(m)}
\]
is only an approximation of $\mathbb E^{\mathbb Q}[H]$. By the central limit theorem,
\[
\frac1M\sum_{m=1}^M H^{(m)}
-
\mathbb E^{\mathbb Q}[H]
=
O_{\mathbb P}\left(\frac1{\sqrt M}\right).
\]

Second, there is a \emph{time-discretization error}. This comes from approximating the continuous-time
process, ordinary integrals, stochastic integrals, or running maxima on a finite grid. This error decreases
as
\[
\Delta t=\frac{T}{N}\to0.
\]

In practice, increasing $M$ reduces statistical error, while increasing $N$ reduces time-discretization error.
\end{solutionpart}
\end{solution}

\end{document}